% Generated by tables/make_methods-compared.py from data/derived/benchmark/methods-compared.json; do not edit.
\subsubsection{A behavioural assumption changed at a limit breach}
\label{sec:results-methods-compared-q9}

This question is answered on case C14. Planted histories of 12 quarters in paired arms, 4,000 per arm: the same breach chain and the same unrelated changes in both arms, the parameter lengthening at a breach with probability one half in the second arm. The two shift tests are calibrated at 5\% on separate null histories; every rule is scored on the same histories. The committed bank documents show that the test operates on one bank, not that the bank behaved as the test flags, so they are not scored here. The run ``Planted histories, paired arms'' reads as follows: 12 quarters, 4,000 histories per arm; its random seed is 7101. Table~\ref{tab:bm-q9-methods} lists the seven methods the document compares, one long-established, one proposed in this work and five reference, with the section of this report that describes each. Tables~\ref{tab:bm-q9-1}, \ref{tab:bm-q9-2}, \ref{tab:bm-q9-3} give the values, one table for each group of metrics that share a run and a layout.

{\footnotesize\setlength{\tabcolsep}{3pt}
% bm-table-methods Q9
\begin{longtable}{>{\raggedright\arraybackslash}p{0.38\textwidth}>{\raggedright\arraybackslash}p{0.17\textwidth}>{\raggedright\arraybackslash}p{0.27\textwidth}>{\raggedright\arraybackslash}p{0.10\textwidth}}
\caption{The methods compared on the question A behavioural assumption changed at a limit breach: the standing of each method, the source and the year the document records for it, and the section of this report that describes it; a dash marks a method with no source recorded or with no section of its own.}\label{tab:bm-q9-methods}\\
\toprule Method & Standing & Source & Section \\ \midrule \endfirsthead
\toprule Method & Standing & Source & Section \\ \midrule \endhead
Any change of the parameter (the plain stability reading) & Long-established & -- & \ref{sec:m22} \\
Breach test, as the engine reads a stated parameter & Proposed in this work & -- & \ref{sec:bs8} \\
Fisher exact test on the flag counts & Reference & -- & -- \\
Mann-Whitney U test on the parameter & Reference & -- & -- \\
Break in level at an unknown date (conformal, 5\%) & Reference & -- & -- \\
Break in level restricted to breach dates (conformal, 5\%) & Reference & -- & -- \\
Welch t test on the parameter & Reference & -- & -- \\
\bottomrule
\end{longtable}
}

{\footnotesize\setlength{\tabcolsep}{3pt}
% bm-table Q9 normal cols=false_flag/lam-0.05;false_flag/lam-0.10;false_flag/lam-0.20
\begin{longtable}{>{\raggedright\arraybackslash}p{0.25\textwidth}>{\raggedleft\arraybackslash}p{0.14\textwidth}>{\raggedleft\arraybackslash}p{0.14\textwidth}>{\raggedleft\arraybackslash}p{0.14\textwidth}}
\caption{A behavioural assumption changed at a limit breach: Flag rate with unrelated changes only (false flags) (share of histories; lower is better), for each method, on the run Planted histories, paired arms. The best value of a column by the declared rule is in bold, ties all in bold; a dash marks a cell without a value, and the notes below the table state why.}\label{tab:bm-q9-1}\\
\toprule Method & Unrelated changes 0.05, breach prevalence 0.10 & Unrelated changes 0.10, breach prevalence 0.25 & Unrelated changes 0.20, breach prevalence 0.50 \\ \midrule \endfirsthead
\toprule Method & Unrelated changes 0.05, breach prevalence 0.10 & Unrelated changes 0.10, breach prevalence 0.25 & Unrelated changes 0.20, breach prevalence 0.50 \\ \midrule \endhead
\multicolumn{4}{l}{\emph{Long-established}} \\
Any change of the parameter (the plain stability reading) & 0.435 & 0.688 & 0.912 \\ % bm-row any_change
\multicolumn{4}{l}{\emph{Proposed in this work}} \\
Breach test, as the engine reads a stated parameter & 0.170 & 0.348 & 0.664 \\ % bm-row engine
\multicolumn{4}{l}{\emph{Reference}} \\
Fisher exact test on the flag counts & \textbf{0.000} & \textbf{0.0013} & \textbf{0.0023} \\ % bm-row fisher
Mann-Whitney U test on the parameter & 0.021 & 0.026 & 0.020 \\ % bm-row mannwhitney
Break in level at an unknown date (conformal, 5\%) & 0.046 & 0.040 & 0.048 \\ % bm-row shift
Break in level restricted to breach dates (conformal, 5\%) & 0.060 & 0.041 & 0.045 \\ % bm-row shift_at_breach
Welch t test on the parameter & \ensuremath{2.5\times10^{-4}} & 0.0033 & 0.013 \\ % bm-row welch
\bottomrule
\end{longtable}
}

{\footnotesize\setlength{\tabcolsep}{3pt}
% bm-table Q9 normal cols=detect/lam-0.05;detect/lam-0.10;detect/lam-0.20
\begin{longtable}{>{\raggedright\arraybackslash}p{0.25\textwidth}>{\raggedleft\arraybackslash}p{0.14\textwidth}>{\raggedleft\arraybackslash}p{0.14\textwidth}>{\raggedleft\arraybackslash}p{0.14\textwidth}}
\caption{A behavioural assumption changed at a limit breach: Flag rate with a change at a breach (detection) (share of histories; higher is better), for each method, on the run Planted histories, paired arms. The best value of a column by the declared rule is in bold, ties all in bold; a dash marks a cell without a value, and the notes below the table state why.}\label{tab:bm-q9-2}\\
\toprule Method & Unrelated changes 0.05, breach prevalence 0.10 & Unrelated changes 0.10, breach prevalence 0.25 & Unrelated changes 0.20, breach prevalence 0.50 \\ \midrule \endfirsthead
\toprule Method & Unrelated changes 0.05, breach prevalence 0.10 & Unrelated changes 0.10, breach prevalence 0.25 & Unrelated changes 0.20, breach prevalence 0.50 \\ \midrule \endhead
\multicolumn{4}{l}{\emph{Long-established}} \\
Any change of the parameter (the plain stability reading) & \textbf{1.000} & \textbf{1.000} & \textbf{1.000} \\ % bm-row any_change
\multicolumn{4}{l}{\emph{Proposed in this work}} \\
Breach test, as the engine reads a stated parameter & \textbf{1.000} & \textbf{1.000} & \textbf{1.000} \\ % bm-row engine
\multicolumn{4}{l}{\emph{Reference}} \\
Fisher exact test on the flag counts & 0.135 & 0.132 & 0.102 \\ % bm-row fisher
Mann-Whitney U test on the parameter & 0.574 & 0.481 & 0.318 \\ % bm-row mannwhitney
Break in level at an unknown date (conformal, 5\%) & 0.122 & 0.058 & 0.013 \\ % bm-row shift
Break in level restricted to breach dates (conformal, 5\%) & 0.685 & 0.068 & 0.048 \\ % bm-row shift_at_breach
Welch t test on the parameter & 0.226 & 0.268 & 0.314 \\ % bm-row welch
\bottomrule
\end{longtable}
}

{\footnotesize\setlength{\tabcolsep}{3pt}
% bm-table Q9 normal cols=auc/lam-0.05;auc/lam-0.10;auc/lam-0.20
\begin{longtable}{>{\raggedright\arraybackslash}p{0.25\textwidth}>{\raggedleft\arraybackslash}p{0.14\textwidth}>{\raggedleft\arraybackslash}p{0.14\textwidth}>{\raggedleft\arraybackslash}p{0.14\textwidth}}
\caption{A behavioural assumption changed at a limit breach: ROC area of the statistic (area; higher is better), for each method, on the run Planted histories, paired arms. The best value of a column by the declared rule is in bold, ties all in bold; a dash marks a cell without a value, and the notes below the table state why.}\label{tab:bm-q9-3}\\
\toprule Method & Unrelated changes 0.05, breach prevalence 0.10 & Unrelated changes 0.10, breach prevalence 0.25 & Unrelated changes 0.20, breach prevalence 0.50 \\ \midrule \endfirsthead
\toprule Method & Unrelated changes 0.05, breach prevalence 0.10 & Unrelated changes 0.10, breach prevalence 0.25 & Unrelated changes 0.20, breach prevalence 0.50 \\ \midrule \endhead
\multicolumn{4}{l}{\emph{Long-established}} \\
Any change of the parameter (the plain stability reading) & 0.901 & 0.863 & 0.848 \\ % bm-row any_change
\multicolumn{4}{l}{\emph{Proposed in this work}} \\
Breach test, as the engine reads a stated parameter & 0.953 & 0.908 & \textbf{0.856} \\ % bm-row engine
\multicolumn{4}{l}{\emph{Reference}} \\
Fisher exact test on the flag counts & \textbf{0.955} & \textbf{0.922} & 0.840 \\ % bm-row fisher
Mann-Whitney U test on the parameter & 0.948 & 0.914 & 0.849 \\ % bm-row mannwhitney
Break in level at an unknown date (conformal, 5\%) & 0.861 & 0.796 & 0.768 \\ % bm-row shift
Break in level restricted to breach dates (conformal, 5\%) & 0.843 & 0.824 & 0.825 \\ % bm-row shift_at_breach
Welch t test on the parameter & 0.776 & 0.761 & 0.774 \\ % bm-row welch
\bottomrule
\end{longtable}
}

The document reads the tables of this question by the declared rule of each metric. Flag rate with unrelated changes only (false flags) (lower is better): best in Fisher exact test on the flag counts 3 of 3. Flag rate with a change at a breach (detection) (higher is better): best in Any change of the parameter (the plain stability reading) 3 of 3, Breach test, as the engine reads a stated parameter 3 of 3; ties count for each method tied. ROC area of the statistic (higher is better): best in Fisher exact test on the flag counts 2 of 3, Breach test, as the engine reads a stated parameter 1 of 3. Breach test, as the engine reads a stated parameter, proposed in this work, is best in 4 of 9 scored cells and not best in 5. A null cell is never ranked; each states why it is null.

The breach test as the engine reads a stated parameter flags unrelated changes far above the nominal 5\%, where the calibrated break-in-level tests stay near it; the table shows it losing on false flags. The product's own manuscript on this method reports it~\citep{proposed:contraste-breach}.
